% ============================================================================
% slides/02-fundamentacao.tex — Capítulo 2: Fundamentação Teórica
% ============================================================================

\section{Fundamentação Teórica}

% --- Algoritmo Memético (AG + 2-opt) ----------------------------------------
\begin{frame}{Algoritmo Memético (AG + $2$-opt)}
  \begin{columns}[c,onlytextwidth]
    \begin{column}{0.53\textwidth}
      \begin{itemize}
        \small
        \item \textbf{Exploração global (AG):} $n_{\text{pop}} = 256$ rotas; seleção por torneio ($\tau = 5$), cruzamento OX, mutação \emph{swap} e elitismo $(\mu + \lambda)$ \autocite{goldberg1989genetic,larranaga1999genetic}.
        \item \textbf{Busca local $2$-opt (\emph{best-improvement}):} remove $(v_i,v_{i+1})$ e $(v_j,v_{j+1})$, reconecta e reverte o segmento intermediário \autocite{croes1958method}:
              \[
                \Delta C = (c_{i,j} + c_{i+1,j+1}) - (c_{i,i+1} + c_{j,j+1}) < 0
              \]
        \item \textbf{Lista estática de candidatos ($k_{\text{viz}} = \kviz$):} reduz a varredura de $\Theta(n^2)$ para $\Theta(n\,k_{\text{viz}})$ sem alterar a função objetivo \autocite{lin1973efficient}, limitada a $\kappa = \maxpassadas$ varreduras.
      \end{itemize}
    \end{column}
    \begin{column}{0.44\textwidth}
      \slideTikz[\slideTallFigH]{%
        \begin{tikzpicture}[node distance=3.4mm]
          \node[sl_term, text width=4.6cm] (init) {Inicialização: $P(0)$ ($n_{\text{pop}}=256$)};
          \node[sl_cpu, below=of init, text width=4.6cm] (gen) {Seleção $\to$ Cruzamento OX $\to$ Mutação \emph{Swap}};
          \node[sl_mem, below=of gen, text width=4.6cm] (ls) {\textbf{Busca Local $2$-opt}\\($k_{\text{viz}}=\kviz$, até $\kappa=\maxpassadas$ varreduras)};
          \node[sl_cpu, below=of ls, text width=4.6cm] (eval) {Avaliação + Elitismo $(\mu+\lambda)$};
          \node[sl_dec, below=4mm of eval, text width=2.3cm] (stop) {Parada atingida?};
          \node[sl_good, below=4mm of stop, text width=3.6cm] (end) {Retorna melhor circuito};

          \draw[sl_arr] (init) -- (gen);
          \draw[sl_arr] (gen) -- (ls);
          \draw[sl_arr] (ls) -- (eval);
          \draw[sl_arr] (eval) -- (stop);
          \draw[sl_arr] (stop) -- node[sl_lbl, right=2pt] {Sim} (end);
          \draw[sl_arr] (stop.east) -- ++(5.5mm,0) |- node[sl_lbl, pos=0.25, right=1pt] {Não ($t{+}1$)} (gen.east);
        \end{tikzpicture}%
      }{Elaborado pelo autor.}
    \end{column}
  \end{columns}
\end{frame}

% --- Contraste Arquitetural CPU vs. GPU -------------------------------------
\begin{frame}{Contraste Arquitetural: CPU vs.\ GPU \autocite{flynn1972taxonomy}}
  \begin{columns}[T,onlytextwidth]
    \begin{column}{0.48\textwidth}
      \slideFigure[\slideDualFigH]{\linewidth}{cpu-gpu-silicon-allocation.png}{Adaptado de \textcite{paz2011gpucpuimage}.}
      \vspace{0.1cm}
      \begin{block}{CPU (MIMD / SIMD)}
        Poucos núcleos complexos, amplo controle e \emph{caches} grandes — otimizada para \textbf{baixa latência} em fluxos sequenciais.
      \end{block}
    \end{column}
    \begin{column}{0.49\textwidth}
      \slideFigure[\slideDualFigH]{\linewidth}{gpu-internal-architecture.png}{Adaptado de \textcite{shah2023gpuarchitectureimage}.}
      \vspace{0.1cm}
      \begin{block}{GPU (SIMT / CUDA)}
        Centenas de ALUs compactas e memória global de alta vazão — otimizada para \textbf{\emph{throughput}} paralelo em dados.
      \end{block}
    \end{column}
  \end{columns}
\end{frame}

% --- Modelo SIMT e Lei de Amdahl --------------------------------------------
\begin{frame}{Modelo SIMT e Gargalo Sequencial (Lei de Amdahl)}
  \begin{columns}[T,onlytextwidth]
    \begin{column}{0.52\textwidth}
      \begin{block}{Execução em GPU (Modelo SIMT) \autocite{nvidia2024cuda}}
        \begin{itemize}
          \item \emph{Threads} escalonadas em \textbf{\emph{warps} de $32$ unidades} executando a mesma instrução.
          \item \textbf{Divergência de fluxo:} desvios condicionais (\texttt{if/else}) irregulares serializam o \emph{warp}.
          \item \textbf{Coalescimento:} leituras contíguas em memória global evitam múltiplas transações.
          \item \textbf{Memória compartilhada:} reduz latência dentro de cada SM (\emph{Streaming Multiprocessor}).
        \end{itemize}
      \end{block}
    \end{column}
    \begin{column}{0.45\textwidth}
      \begin{alertblock}{Lei de Amdahl \autocite{amdahl1967validity}}
        \[
          S = \frac{1}{(1 - f) + \dfrac{f}{p}} < \frac{1}{1 - f}
        \]
        \begin{itemize}
          \item $f$: fração paralelizável; $p$: aceleração obtida.
          \item Ao acelerar intensamente apenas a busca local ($p \gg 1$), a fração restante $(1 - f)$ em CPU passa a dominar o tempo total.
        \end{itemize}
      \end{alertblock}
    \end{column}
  \end{columns}
\end{frame}

